% Reviewed translation: PDF pages 56--60 / printed pages 42--46.
% Continues the Proof of Lemma 3.2 from the immutable approved sample derivative.

\MathBookSourcePage{56}
这里 $q$ 是下列 Neumann 问题的唯一解(相差一个加法常数):
\[
\left\{\begin{aligned}
\Delta q&=0&&\text{in }\Omega,\\
\partial q/\partial n&=\mathbf g\cdot\mathbf n&&\text{on }\Gamma.
\end{aligned}\right.
\]
而 $\phi$ 是下列双调和方程的唯一解:
\[
\Delta^2\phi=0\quad\text{in }\Omega,
\]
\[
\phi=0\quad\text{on }\Gamma,
\qquad
\partial\phi/\partial n=\partial q/\partial\tau-\mathbf g\cdot\boldsymbol\tau
\quad\text{on }\Gamma.
\]
读者容易验证 $\mathbf u$ 满足该 Lemma 的结论.
\end{Proof}

现在转向空间 $V$. 根据 Theorem~\ref{thm:sec3-stream-function-characterization},
每个函数 $\mathbf v\in V$ 都有流函数 $\phi\in H^1(\Omega)$, 即
\[
\mathbf v=\Curl\phi.
\]
同样, 如果规定 $\phi=0$ on $\Gamma_0$, 则 $\phi$ 被唯一确定. 在这种情形下容易验证
$\phi$ 属于空间
\begin{equation}\label{eq:sec3-psi-space}
\Psi=\{\chi\in H^2(\Omega);\ \chi|_{\Gamma_0}=0,\
\ \chi|_{\Gamma_i}\text{ is constant for }1\le i\le p,\
\ \partial\chi/\partial n=0\text{ on }\Gamma\}.
\end{equation}
注意, 半范数 $\|\Delta\phi\|_{0,\Omega}$ 是 $\Psi$ 上的范数, 且与
$\|\phi\|_{2,\Omega}$ 等价. 事实上, 当 $\phi\in\Psi$ 时,
$\Grad\phi\in H_0^1(\Omega)^2$; 仿照 Section~\ref{subsec:biharmonic-dirichlet-example}
并利用 Lemma~\ref{lem:sec3-boundary-poincare}, 容易证明
\begin{equation}\label{eq:sec3-psi-laplacian-norm}
\|\Delta\phi\|_{0,\Omega}=|\phi|_{2,\Omega}
\ge C_1|\phi|_{1,\Omega}\ge C_2\|\phi\|_{0,\Omega}.
\end{equation}
利用这个空间, 可以对 $V$ 给出 Corollary~\ref{cor:sec3-h-stream-functions} 和
Theorem~\ref{thm:sec3-orthogonal-decomposition} 的类似结论.

\setcounter{statement}{1}
\begin{Corollary}\label{cor:sec3-v-stream-functions}
空间 $V$ 满足恒等式
\[
V=\{\Curl\phi;\ \phi\in\Psi\},
\]
并且映射 $\Curl$ 是从 $\Psi$ 到 $V$ 的同构. 此外, $\mathbf v$ 的流函数 $\phi$
可刻画为下列问题的唯一解:
\begin{equation}\label{eq:sec3-v-stream-variational}
\phi\in\Psi,\qquad
(\Delta\phi,\Delta\chi)=-(\Curl\mathbf v,\Delta\chi)
\quad\forall\chi\in\Psi.
\end{equation}
\end{Corollary}

\setcounter{statement}{2}
\begin{Remark}\label{rem:sec3-v-stream-biharmonic-interpretation}
问题 \eqref{eq:sec3-v-stream-variational} 有如下解释:
\begin{equation}\label{eq:sec3-v-stream-biharmonic-system}
\left\{\begin{aligned}
\Delta^2\phi&=-\Delta(\Curl\mathbf v)&&\text{in }H^{-2}(\Omega),\\
\phi|_{\Gamma_0}&=0,\qquad
\phi|_{\Gamma_i}=c_i&&1\le i\le p,\\
\partial\phi/\partial n&=0&&\text{on }\Gamma,\\
\int_{\Gamma_i}\partial(\Delta\phi+\Curl\mathbf v)/\partial n\dd s&=0
&&1\le i\le p.
\end{aligned}\right.
\end{equation}

\MathBookSourcePage{57}
事实上, 把 $\chi$ 限制在 $\D(\Omega)$ 中, 立即得到
\eqref{eq:sec3-v-stream-biharmonic-system} 的第一个等式; 将该等式与
$\chi\in\Psi$ 作标量积, 分部积分并与 \eqref{eq:sec3-v-stream-variational} 比较, 形式上得到
\[
\langle\partial(\Delta\phi+\Curl\mathbf v)/\partial n,\chi\rangle_\Gamma=0
\quad\forall\chi\in\Psi.
\]
由于 $\chi\in\Psi$, 这又形式上蕴含
\[
\int_{\Gamma_i}\partial(\Delta\phi+\Curl\mathbf v)/\partial n\dd s=0
\quad\text{for }1\le i\le p.
\]
\end{Remark}

\setcounter{statement}{2}
\begin{Theorem}\label{thm:sec3-h01-orthogonal-decomposition}
$H_0^1(\Omega)^2$ 中每个函数 $\mathbf v$ 都有如下正交分解:
\begin{equation}\label{eq:sec3-h01-orthogonal-decomposition}
\mathbf v=(-\Delta)^{-1}\Grad q+\Curl\phi
\qquad\text{(see Definition~\ref{def:sec2-greens-operator})},
\end{equation}
其中 $\phi\in\Psi$ 是问题 \eqref{eq:sec3-v-stream-variational} 的唯一解,
而 $q$ 由 \eqref{eq:sec3-h01-orthogonal-decomposition} 在 $L_0^2(\Omega)$ 中唯一确定.
\end{Theorem}
\begin{Proof}
分解 \eqref{eq:sec3-h01-orthogonal-decomposition} 立即由
Corollary~\ref{cor:sec2-vperp-green} 和
Corollary~\ref{cor:sec3-v-stream-functions} 得到. 因此只需检验 $\phi$ 满足
\eqref{eq:sec3-v-stream-variational}. 由 \eqref{eq:sec3-h01-orthogonal-decomposition} 推得
\[
\Curl\mathbf v=-\Delta\phi+\Curl(-\Delta)^{-1}\Grad q.
\]
对 $\mathbf w\in\mathcal V$ 用 $\Curl\mathbf w$ 乘这个等式的两边, 得
\begin{align*}
(\Curl\mathbf v,\Curl\mathbf w)
&=(-\Delta\phi,\Curl\mathbf w)
+(\Curl(-\Delta)^{-1}\Grad q,\Curl\mathbf w)\\
&=(-\Delta\phi,\Curl\mathbf w)
+((-\Delta)^{-1}\Grad q,\Curl(\Curl\mathbf w))\\
&=(-\Delta\phi,\Curl\mathbf w)
+((-\Delta)^{-1}\Grad q,-\Delta\mathbf w),
\end{align*}
因为 $\Div\mathbf w=0$. 因此
\begin{align*}
(\Curl\mathbf v,\Curl\mathbf w)
&=(-\Delta\phi,\Curl\mathbf w)+(\Grad q,\mathbf w)\\
&=(-\Delta\phi,\Curl\mathbf w)
\qquad\forall\mathbf w\in\mathcal V.
\end{align*}
由于 $\mathcal V$ 在 $V$ 中稠密, 这意味着
\[
(-\Delta\phi,\Curl\mathbf w)=(\Curl\mathbf v,\Curl\mathbf w)
\quad\forall\mathbf w\in V,
\]
而这确实等价于 \eqref{eq:sec3-v-stream-variational}.
\end{Proof}

\setcounter{subsection}{1}
\subsection{$H(\Div;\Omega)\cap H(\Curl;\Omega)$ 中函数的正则性应用}
\label{subsec:sec3-hdiv-hcurl-regularity}

在 Lemma~\ref{lem:sec2-h1-hdiv-hcurl} 中已经看到,
$H_0(\Div;\Omega)\cap H_0(\Curl;\Omega)$ 中的函数属于 $H_0^1(\Omega)^N$,
而 $H(\Div;\Omega)\cap H(\Curl;\Omega)$ 中的函数局部属于 $H^1(\Omega)^N$.
不过, 在实际中常会遇到只满足其中一个边界条件的中间情形, 此时自然产生的问题是:
在边界 $\Gamma$ 附近会发生什么? 如果 $\Gamma$ 足够光滑, Duvaut \& Lions [26]
和 Gobert [40] 给出的答案是: 如果 $\mathbf u\cdot\mathbf n=0$ 或
$\mathbf u\times\mathbf n=\mathbf0$, 则 $\mathbf u\in H^1(\Omega)^N$.

\MathBookSourcePage{58}
然而, 当 $\Gamma$ 有角点时, 他们的方法并不适用. 事实表明, 只要 $\Omega$ 没有凹角
(至少在二维情形中), 这一性质仍然成立; 而在二维情形中, Subsection~\ref{subsec:two-dimensional-decomposition}
的结果给出一个非常简单的证明.

回顾 $\Omega$ 是 $\mathbb R^2$ 中连通、有界的开子集. 置
\[
U=H_0(\Div;\Omega)\cap H(\Curl;\Omega),
\qquad
W=H(\Div;\Omega)\cap H_0(\Curl;\Omega).
\]

\setcounter{statement}{0}
\begin{Proposition}\label{prop:sec3-hdiv-hcurl-h1-inclusion}
如果 $\Gamma$ 属于 $\mathcal C^{1,1}$ 类, 或者 $\Gamma$ 分片光滑且没有凹角,
则 $U$ 和 $W$ 在代数意义和拓扑意义下都包含于 $H^1(\Omega)^2$ 中.
\end{Proposition}
\begin{Proof}
$1^\circ)$ 设 $\mathbf v\in U$. 根据 Theorem~\ref{thm:sec3-orthogonal-decomposition}
以及 Remark~\ref{rem:sec3-stream-boundary-problem} 和
Remark~\ref{rem:sec3-neumann-interpretation}, $\mathbf v$ 可分解为
\begin{equation*}
\tag{\ref*{eq:sec3-orthogonal-decomposition}}
\mathbf v=\Grad q+\Curl\phi,
\end{equation*}
其中 $q$ 是下列齐次 Neumann 问题的解(相差一个加法常数唯一):
\begin{equation}\label{eq:sec3-regularity-neumann}
\left\{\begin{aligned}
\Delta q&=\Div\mathbf v&&\text{in }\Omega,\\
\partial q/\partial n&=0&&\text{on }\Gamma,
\end{aligned}\right.
\end{equation}
而 $\phi$ 是下列非齐次 Dirichlet 问题的解:
\[
-\Delta\phi=\Curl\mathbf v\quad\text{in }\Omega,
\]
\[
\phi|_{\Gamma_0}=0,
\qquad
\phi|_{\Gamma_i}=c_i,
\]
其中常数 $c_i$ 由 \eqref{eq:sec3-stream-variational} 确定.
由于 $\Div\mathbf v$ 和 $\Curl\mathbf v$ 都属于 $L^2(\Omega)$, 并且 $\Gamma$
满足上述假设, 由 Theorem~\ref{thm:laplace-dirichlet-regularity} 和
Theorem~\ref{thm:laplace-neumann-regularity}, $q$ 和 $\phi$ 都属于 $H^2(\Omega)$, 且
\begin{equation}\label{eq:sec3-regularity-potential-estimates}
\left\{\begin{aligned}
\|q\|_{2,\Omega}&\le C_1\|\Div\mathbf v\|_{0,\Omega},\\
\|\phi\|_{2,\Omega}&\le C_2(\|\Curl\mathbf v\|_{0,\Omega}
+\|\gamma_0\phi\|_{3/2,\Gamma}).
\end{aligned}\right.
\end{equation}
因为 $\phi$ 在每个 $\Gamma_i$ 上为常数且在 $\Gamma_0$ 上为零, 有
\[
\|\gamma_0\phi\|_{3/2,\Gamma}
\le C_3\|\phi\|_{1,\Omega}\le C_4|\phi|_{1,\Omega}.
\]
此外, 由 \eqref{eq:sec3-stream-variational} 和
\eqref{eq:sec3-curl-remainder-variational}, 得
\[
|\phi|_{1,\Omega}\le\|\mathbf v-\Grad q\|_{0,\Omega}
\le\|\mathbf v\|_{0,\Omega}.
\]
因此, 汇总这些不等式并代入 \eqref{eq:sec3-orthogonal-decomposition}, 推得
$\mathbf v\in H^1(\Omega)^2$, 且
\begin{equation}\label{eq:sec3-hdiv-hcurl-h1-bound}
\|\mathbf v\|_{1,\Omega}\le C_5(\|\mathbf v\|_{0,\Omega}
+\|\Div\mathbf v\|_{0,\Omega}+\|\Curl\mathbf v\|_{0,\Omega}).
\end{equation}

$2^\circ)$ 对空间 $W$, 可以使用 Remark~\ref{rem:sec2-hcurl-hdiv-2d} 的简单技巧:
对 $W$ 中的 $\mathbf u$, 向量 $\mathbf v=(-u_2,u_1)$ 属于 $U$, 且
$\Div\mathbf v=-\Curl\mathbf u$, $\Curl\mathbf v=\Div\mathbf u$,
$\mathbf v\cdot\mathbf n=-\mathbf u\cdot\boldsymbol\tau$.
因此由 $1^\circ)$ 推得 $\mathbf v\in H^1(\Omega)^2$ 并满足估计
\eqref{eq:sec3-hdiv-hcurl-h1-bound}; $\mathbf u$ 也具有相同性质.
\end{Proof}

\MathBookSourcePage{59}
这个 Proposition 有一些有意义的推论, 还需要作若干说明.

\setcounter{statement}{3}
\begin{Remark}\label{rem:sec3-reentrant-corner-counterexample}
存在带凹角的多边形区域 $\Omega$, 使问题 \eqref{eq:sec3-regularity-neumann}
的解 $q$ 不属于 $H^2(\Omega)$(参见 Grisvard [42]). 因此, 在
\eqref{eq:sec3-orthogonal-decomposition} 中取 $\phi=0$, 就得到一个属于 $U$ 但不属于
$H^1(\Omega)^2$ 的函数 $\mathbf v$. 所以, 一般而言, 对多边形区域角度的条件既必要又充分.
\end{Remark}

\setcounter{statement}{4}
\begin{Remark}\label{rem:sec3-simply-connected-norm-equivalence}
设 $\Omega$ 单连通且满足 Proposition~\ref{prop:sec3-hdiv-hcurl-h1-inclusion} 的条件.
由于 $p=0$, 由 \eqref{eq:sec3-regularity-potential-estimates} 推得 $U$ 中有范数等价
\begin{equation}\label{eq:sec3-simply-connected-h1-equivalence}
\|\mathbf v\|_{1,\Omega}\cong
\|\Div\mathbf v\|_{0,\Omega}+\|\Curl\mathbf v\|_{0,\Omega}.
\end{equation}
这又蕴含 $H\cap\operatorname{Ker}(\Curl)=\{\mathbf0\}$, 如
Remark~\ref{rem:sec2-h-curlzero} 中所述.

另一方面, 如果 $\Omega$ 不是单连通的, 则等价关系
\eqref{eq:sec3-simply-connected-h1-equivalence} 不再成立, 因为
$H\cap\operatorname{Ker}(\Curl)$ 中存在非零函数.
\end{Remark}

\setcounter{statement}{5}
\begin{Remark}\label{rem:sec3-extension-hdiv-hcurl}
设 $\Omega$ 满足 Proposition~\ref{prop:sec3-hdiv-hcurl-h1-inclusion} 的条件.
令 $P\in\mathcal L(H^1(\Omega);H^1(\mathbb R^2))$ 为
Theorem~\ref{thm:sobolev-density-extension} 的延拓算子. 将 $P$ 逐分量应用于 $U$(或 $W$)
中的函数, 也给出从 $U$(或 $W$) 到 $H^1(\mathbb R^2)^2$ 的延拓.
\end{Remark}

\setcounter{subsection}{2}
\subsection{三维向量场的分解}\label{subsec:three-dimensional-decomposition}

\setcounter{statement}{3}
\begin{Theorem}\label{thm:sec3-three-dimensional-vector-potential}
向量场 $\mathbf v\in L^2(\Omega)^3$ 满足条件
\eqref{eq:sec3-divfree-zero-flux}, 当且仅当存在向量势
$\boldsymbol\phi\in H^1(\Omega)^3$, 使得
\[
\mathbf v=\Curl\boldsymbol\phi.
\]
此外,
\[
\Div\boldsymbol\phi=0.
\]
\end{Theorem}
\begin{Proof}
证明与 Theorem~\ref{thm:sec3-stream-function-characterization} 的证明十分相似.

$1^\circ)$ 设 $\boldsymbol\phi\in H^1(\Omega)^3$ 且
$\mathbf v=\Curl\boldsymbol\phi$. 则 $\Div\mathbf v=0$, 还需验证
$\langle\mathbf v\cdot\mathbf n,1\rangle_{\Gamma_i}=0$.
对 $0\le i\le p$, 取 $\D(\mathbb R^3)$ 中的函数 $\theta_i$, 满足
\[
0\le\theta_i(x)\le1\quad\text{in }\mathbb R^3,
\qquad
\theta_i(x)=\delta_{ij}\quad\text{in a neighborhood of }\Gamma_j.
\]
置 $\mathbf w_i=\Curl(\theta_i\boldsymbol\phi)$. 显然
$\mathbf w_i\in L^2(\Omega)^3$ 且 $\Div\mathbf w_i=0$. 此外,
\[
\mathbf w_i\cdot\mathbf n|_{\Gamma_j}=0\quad\text{if }j\ne i,
\qquad
\mathbf w_i\cdot\mathbf n|_{\Gamma_i}
=\mathbf v\cdot\mathbf n|_{\Gamma_i}.
\]
因此
\[
\langle\mathbf v\cdot\mathbf n,1\rangle_{\Gamma_i}
=\langle\mathbf w_i\cdot\mathbf n,1\rangle_\Gamma
=\int_\Omega\Div\mathbf w_i\dd x=0
\quad\text{for }0\le i\le p.
\]

\MathBookSourcePage{60}
$2^\circ)$ 反之, 设 $\mathbf v\in L^2(\Omega)^3$ 满足
\eqref{eq:sec3-divfree-zero-flux}. 与二维情形相同, 将 $\mathbf v$ 延拓到整个空间,
使延拓函数 $\widetilde{\mathbf v}$ 属于 $L^2(\mathbb R^3)^3$, 无散且具有紧支集.
因为 $\widetilde{\mathbf v}$ 的支集紧, 它的 Fourier 变换
$\mathscr F\widetilde{\mathbf v}$ 在 $\mathbb R^3$ 中是解析的. 用 Fourier 变换表示,
$\Div\widetilde{\mathbf v}=0$ 和
$\widetilde{\mathbf v}=\Curl\boldsymbol\phi$ 分别变为
\begin{equation}\label{eq:sec3-three-d-fourier-divfree}
\sum_{i=1}^3\mu_i\mathscr F\widetilde v_i=0,
\end{equation}
以及
\begin{equation}\label{eq:sec3-three-d-fourier-curl}
\mathscr F\widetilde{\mathbf v}=2i\pi(
\mu_2\mathscr F\phi_3-\mu_3\mathscr F\phi_2,
\mu_3\mathscr F\phi_1-\mu_1\mathscr F\phi_3,
\mu_1\mathscr F\phi_2-\mu_2\mathscr F\phi_1).
\end{equation}
注意, 如果 \eqref{eq:sec3-three-d-fourier-divfree} 成立, 则
\eqref{eq:sec3-three-d-fourier-curl} 与条件
\begin{equation}\label{eq:sec3-three-d-fourier-potential-divfree}
\sum_{i=1}^3\mu_i\mathscr F\phi_i=0
\end{equation}
唯一确定 $\mathscr F\boldsymbol\phi$. 方程
\eqref{eq:sec3-three-d-fourier-divfree}--\eqref{eq:sec3-three-d-fourier-potential-divfree}
的唯一解是
\[
\mathscr F\boldsymbol\phi=\frac{1}{2i\pi\|\boldsymbol\mu\|^2}(
\mu_3\mathscr F\widetilde v_2-\mu_2\mathscr F\widetilde v_3,
\mu_1\mathscr F\widetilde v_3-\mu_3\mathscr F\widetilde v_1,
\mu_2\mathscr F\widetilde v_1-\mu_1\mathscr F\widetilde v_2).
\]
显然, 这些方程蕴含 $\mu_j\mathscr F\phi_i\in L^2(\mathbb R^3)$, 并且
\[
|\mathscr F\phi_i|\le
(|\mathscr F\widetilde v_j|+|\mathscr F\widetilde v_k|)/(2\pi\|\boldsymbol\mu\|).
\]
因此, 只要 $\mathscr F\boldsymbol\phi$ 在原点处有界, 它的逆变换就属于
$H^1(\Omega)^3$. 首先注意到, \eqref{eq:sec3-three-d-fourier-divfree} 蕴含
$\mathscr F\widetilde{\mathbf v}(0)=\mathbf0$. 又因为
$\mathscr F\widetilde{\mathbf v}$ 解析, 所以在零点的邻域中
\[
\mathscr F\widetilde{\mathbf v}(\boldsymbol\mu)
=\sum_{j=1}^3\mu_j
(\partial\mathscr F\widetilde{\mathbf v}(0)/\partial\mu_j)
+O(\|\boldsymbol\mu\|^2).
\]
所以当 $\boldsymbol\mu\to0$ 时, $\mathscr F\boldsymbol\phi$ 有界.

把 $\mathscr F\boldsymbol\phi$ 的逆变换 $\boldsymbol\phi$ 限制在 $\Omega$ 上,
便得到 $H^1(\Omega)^3$ 中的函数 $\boldsymbol\phi$, 使得
$\mathbf v=\Curl\boldsymbol\phi$, 并且 $\Div\boldsymbol\phi=0$.
\end{Proof}

\setcounter{statement}{6}
\begin{Remark}\label{rem:sec3-divfree-vector-potential-choice}
上述构造表明, 总可以选择无散的向量势. 当然这不是唯一可能的选择. 例如, 可以把
\eqref{eq:sec3-three-d-fourier-potential-divfree} 换成
\begin{equation}\tag{\ref*{eq:sec3-three-d-fourier-potential-divfree}$'$}
\label{eq:sec3-three-d-fourier-potential-third-zero}
\mathscr F\phi_3=0.
\end{equation}
则 \eqref{eq:sec3-three-d-fourier-divfree}, \eqref{eq:sec3-three-d-fourier-curl}
和 \eqref{eq:sec3-three-d-fourier-potential-third-zero} 的唯一解为
\[
\mathscr F\phi_1=\mathscr F\widetilde v_2/(2i\pi\mu_3),
\qquad
\mathscr F\phi_2=-\mathscr F\widetilde v_1/(2i\pi\mu_3),
\qquad
\mathscr F\phi_3=0.
\]
$\mathbf v$ 的相应向量势 $\boldsymbol\phi$ 具有形式 $(\phi_1,\phi_2,0)$.
\end{Remark}

\setcounter{statement}{7}
\begin{Remark}\label{rem:sec3-vector-potential-poisson}
$\mathbf v$ 的每个无散向量势 $\boldsymbol\phi$ 都满足
\[
-\Delta\boldsymbol\phi=\Curl\mathbf v\quad\text{in }H^{-1}(\Omega)^3.
\]
\end{Remark}
